% BHU PYQ -- Manually transcribed & error-corrected
\documentclass[11pt,a4paper]{article}

\usepackage[a4paper,top=2.5cm,bottom=2.5cm,left=2.8cm,right=2.8cm,headheight=14pt]{geometry}
\usepackage{amsmath,amssymb}
\usepackage[T1]{fontenc}
\usepackage[utf8]{inputenc}
\usepackage{lmodern}
\usepackage{microtype}
\usepackage{fancyhdr}
\usepackage{enumitem}
\usepackage{hyperref}
\usepackage{lastpage}
\usepackage{parskip}

\hypersetup{colorlinks=true,urlcolor=black,linkcolor=black,
  pdftitle={BPT-505: Electromagnetic Theory -- BHU PYQ},pdfauthor={Banaras Hindu University}}

\pagestyle{fancy}\fancyhf{}
\renewcommand{\headrulewidth}{0.4pt}\renewcommand{\footrulewidth}{0.4pt}
\fancyhead[L]{\small\textit{Banaras Hindu University}}
\fancyhead[R]{\small\textit{BPT-505: Electromagnetic Theory}}
\fancyfoot[C]{\small Page \thepage\ of \pageref{LastPage}}

\newlist{parts}{enumerate}{2}
\setlist[parts,1]{label=(\alph*),leftmargin=*,itemsep=5pt,topsep=3pt}
\setlist[parts,2]{label=(\roman*),leftmargin=*,itemsep=3pt,topsep=2pt}
\newcommand{\pts}[1]{\hfill\textit{[#1]}}

\begin{document}

\begin{center}
  {\large\bfseries BANARAS HINDU UNIVERSITY}\\[2pt]
  {\normalsize B.Sc. (Hons.) Semester V Examination, 2022-23}\\[6pt]
  \rule{0.8\linewidth}{0.5pt}\\[6pt]
  {\large\bfseries Paper No. BPT-505: Electromagnetic Theory}\\[4pt]
  {\normalsize Time: 3 Hours\hfill Maximum Marks: 70}
\end{center}
\rule{\linewidth}{0.4pt}
\smallskip
\noindent\textit{Instructions: Answer all questions. Symbols have their usual meanings.}
\bigskip

\noindent\textbf{Question 1.} Answer any \textbf{four} of the following: \pts{4 $\times$ 5 = 20}
\begin{parts}
  \item For the vector field:
    \\[ \vec{v} = (2xz + 3y^2) \hat{i} + (4yz^2) \hat{j} \]
    calculate the line integral along a straight-line path from $(0,0,0)$ to $(1,1,1)$.
  \item If $\hat{r}$ is the unit radial vector defined as $\hat{r} = \vec{r}/r$, where $\vec{r} = x\hat{i} + y\hat{j} + z\hat{k}$, compute the gradient expression $(\hat{r} \cdot \nabla)\hat{r}$.
  \item Show that Maxwell's equations in electrodynamics are consistent with the continuity equation.
  \item Find the skin depth of copper with conductivity $\sigma = 5.8 \times 10^7\,\text{mho/m}$ and relative permeability $\mu_r = 1$ at frequency $10^9\,\text{Hz}$.
  \item In a non-magnetic medium, the electric field is given by:
    \\[ \vec{E} = 4 \sin(2\pi \times 10^7 t - 0.8 x) \hat{j} \,\text{V/m} \]
    Compute the time-average power carried by the wave.
\end{parts}

\medskip\rule{\linewidth}{0.2pt}

\noindent\textbf{Question 2.}
\begin{parts}
  \item In a lossy dielectric medium, show that $\tan\delta = \epsilon''/\epsilon'$, where $\delta$ is the loss angle of the medium. Show how the phase of the intrinsic impedance is related to the loss angle. \pts{8}
  \item Consider the multipole expansion of the electrostatic potential and show that the dipole term can be written in the form:
    \\[ V(\vec{r}) = \frac{1}{4\pi\epsilon_0} \frac{\vec{p} \cdot \vec{r}}{r^3} \]
    where $\vec{p}$ is the electric dipole moment. \pts{6}
\end{parts}

\medskip\rule{\linewidth}{0.2pt}

\noindent\textbf{Question 3.}
\begin{parts}
  \item Consider Ampere's law in the presence of magnetic materials. Write it in both differential and integral forms. Obtain the magnetostatic boundary conditions using the integral form of the equations. \pts{8}
  \item Discuss the propagation of electromagnetic waves in a lossy dielectric medium and derive the expression for the intrinsic impedance of the medium. \pts{6}
\end{parts}
\hfill \textbf{OR}
\begin{parts}
  \item What do you mean by the divergence of a vector field? Calculate the divergence of:
    \\[ \vec{v} = y^2 \hat{i} + (2xy + z^2) \hat{j} + 2yz \hat{k} \] \pts{7}
  \item Explain the meaning of a gauge transformation. Write the Lorentz gauge condition and show that both electrostatic and vector potentials obey the wave equation in the Lorentz gauge in the absence of sources. \pts{7}
\end{parts}

\medskip\rule{\linewidth}{0.2pt}

\noindent\textbf{Question 4.}
Obtain the electrostatic potential both inside and outside of a hollow sphere of radius $R$ which is maintained at a constant surface charge density $\sigma_0$. \pts{12}

\medskip\rule{\linewidth}{0.2pt}

\noindent\textbf{Question 5.}
\begin{parts}
  \item Show that the $TE_{10}$ mode is the dominant mode (lowest cutoff frequency) of all the modes supported by a rectangular waveguide. \pts{8}
  \item Derive the relation between reflection and transmission coefficients for an obliquely incident electromagnetic wave at a dielectric interface where the wave electric field is parallel to the plane of incidence (parallel polarization). \pts{6}
\end{parts}
\hfill \textbf{OR}
\begin{parts}
  \item A plane electromagnetic wave is normally incident on the boundary $z = 0$ formed by a lossless dielectric medium and a good conductor. Obtain the $z$-values at which the minimum value of $|E_{1y}|$ occurs in medium 1, where $E_{1y}$ is the total electric field in medium 1. \pts{8}
  \item Show that a $TE_{10}$ mode of a rectangular waveguide can be represented as a sum of two plane TEM waves propagating along zigzag paths between the guide walls. \pts{6}
\end{parts}

\vfill
\rule{\linewidth}{0.4pt}
\begin{center}\small
  \href{https://bhu-exam.vercel.app/}{Click here to visit our website}\\[4pt]
  \href{https://me-aryan.vercel.app/}{Click here to visit our contributor}\\[4pt]
  \href{https://quashan-me.vercel.app/}{Click here to visit our contributor}
\end{center}

\end{document}
